% Docx build preamble.
%
% build-resume.sh splices this in front of Paul_Walko-Resume.tex's \begin{document}
% body, replacing the LaTeX preamble, then runs pandoc. Redefining the resume's
% macros here means the .tex stays the single source of content -- only the
% rendering of each macro changes.
%
% Two things matter for ATS parsers reading a .docx:
%   * No inline math. Pandoc turns $...$ into an OMML equation, which resume
%     parsers cannot read at all. Hence \sep becomes a literal "|".
%   * Job/education entries must be paragraphs, not list items, so the parser
%     sees "title / employer | location | dates" as running text. Only the
%     accomplishment bullets stay bulleted.

\documentclass{article}
\usepackage{hyperref}

\newcommand{\sep}{|}

\newcommand{\resumeSubHeadingListStart}{}
\newcommand{\resumeSubHeadingListEnd}{}

\newcommand{\resumeEducation}[4]{\par\noindent\textbf{#1}\newline #2\newline #3 \sep{} #4\par}
\newcommand{\resumeSubheading}[4]{\par\noindent\textbf{#1}\newline #2\newline #4 \sep{} #3\par}
\newcommand{\resumeProjectHeading}[1]{\par\noindent #1\par}

% Workday drops the list marker on the FIRST item of every Word list, so the
% opening bullet of each role description arrived unmarked. Emitting the bullet
% as a literal character in a plain paragraph puts it in the text run, where no
% list-handling logic can strip it. Skills/Talks below stay real Word lists -- a
% literal bullet there would end up glued onto the first extracted skill.
\newcommand{\resumeItemListStart}{}
\newcommand{\resumeItemListEnd}{}
\newcommand{\resumeItem}[1]{\par\noindent • #1\par}
\newcommand{\resumeSubItem}[1]{\resumeItem{#1}}

\newcommand{\resumeTextListStart}{\begin{itemize}}
\newcommand{\resumeTextListEnd}{\end{itemize}}
\newcommand{\resumeTextItem}[1]{\item #1}
